\section{Fourier Analysis of Boolean Functions}

Let start motivating the topic constructing a polynomial that represents a boolean function.

\begin{example}[Fourier representation of Max boolean function]
    Let $Max_2:\{-1,1\}^2 \rightarrow \{-1,1\}$, which takes the max of the input. 
    We construct the fourier representation in the following way:
    $Max_2(x_1,x_2) = \frac{1}{2} + \frac{1}{2} x_1 + \frac{1}{2} x_2 + \frac{1}{2} x_1x_2.$
\end{example}

\begin{theorem}
    Any function $f:\{+1,-1\}^{n} \rightarrow \mathbb{R}$ is uniquely expressible as a multilinear polynomial.
\end{theorem}

We represent function as $f(x) = f(x_1,...,x_n)= \sum_{S \subseteq [n]} \hat{f}(S) \prod_{i \in S} x_i.$
In this way we can see that the fourier expansion is a linear combinations "XOR"s functions. This is 
because $\prod_{i \in S} x_i$ is the parity or XOR of bits in $S$.


